\documentclass[11pt,a4paper]{article}

\usepackage[spanish,es-nodecimaldot]{babel}
\usepackage[utf8]{inputenc}
\usepackage[T1]{fontenc}
\usepackage{lmodern}
\usepackage{geometry}
\usepackage{setspace}
\usepackage{microtype}
\usepackage{xcolor}
\usepackage{graphicx}
\usepackage{booktabs}
\usepackage{longtable}
\usepackage{tabularx}
\usepackage{array}
\usepackage{enumitem}
\usepackage{float}
\usepackage{amsmath}
\usepackage{hyperref}
\usepackage{fancyhdr}
\usepackage{titlesec}
\usepackage{tcolorbox}
\usepackage{listings}
\usepackage{caption}
\usepackage{lastpage}

\geometry{top=2.0cm,bottom=2.0cm,left=2.2cm,right=2.2cm}
\setstretch{1.08}
\setlength{\parindent}{0pt}
\setlength{\parskip}{0.45em}
\setlist{itemsep=0.25em,topsep=0.35em}

\definecolor{ineblue}{HTML}{123B5D}
\definecolor{inecyan}{HTML}{1E7894}
\definecolor{riskred}{HTML}{C9351E}
\definecolor{riskamber}{HTML}{C47D10}
\definecolor{riskgreen}{HTML}{257E4E}
\definecolor{lightblue}{HTML}{EAF3F7}
\definecolor{lightgray}{HTML}{F3F4F6}
\definecolor{darkgray}{HTML}{4B5563}
\definecolor{codebg}{HTML}{F7F8FA}

\hypersetup{
    colorlinks=true,
    linkcolor=ineblue,
    urlcolor=inecyan,
    pdftitle={Manual de usuario del sistema de riesgo de abstención electoral},
    pdfauthor={Equipo Loop-E}
}

\pagestyle{fancy}
\fancyhf{}
\lhead{\small CCPC--INE \textbar{} Manual de usuario}
\rhead{\small Predicción y escenarios}
\cfoot{\small Página \thepage\ de \pageref{LastPage}}
\setlength{\headheight}{14pt}

\titleformat{\section}{\Large\bfseries\color{ineblue}}{\thesection.}{0.45em}{}
\titleformat{\subsection}{\large\bfseries\color{inecyan}}{\thesubsection.}{0.45em}{}
\titleformat{\subsubsection}{\normalsize\bfseries\color{darkgray}}{\thesubsubsection.}{0.45em}{}

\newcommand{\ruta}[1]{\texttt{\small #1}}
\newcommand{\comando}[1]{\texttt{\small #1}}
\newcommand{\alto}{\textcolor{riskred}{\textbf{alto}}}
\newcommand{\medio}{\textcolor{riskamber}{\textbf{medio}}}
\newcommand{\bajo}{\textcolor{riskgreen}{\textbf{bajo}}}

\newtcolorbox{objetivo}{
    colback=lightblue,colframe=inecyan,boxrule=0.6pt,arc=3pt,
    left=8pt,right=8pt,top=6pt,bottom=6pt
}
\newtcolorbox{importante}{
    colback=yellow!8,colframe=riskamber,boxrule=0.6pt,arc=3pt,
    left=8pt,right=8pt,top=6pt,bottom=6pt
}
\newtcolorbox{resultado}{
    colback=green!5,colframe=riskgreen,boxrule=0.6pt,arc=3pt,
    left=8pt,right=8pt,top=6pt,bottom=6pt
}

\lstdefinestyle{terminal}{
    basicstyle=\ttfamily\footnotesize,
    backgroundcolor=\color{codebg},
    frame=single,
    rulecolor=\color{lightgray},
    breaklines=true,
    columns=fullflexible,
    keepspaces=true,
    showstringspaces=false,
    language=bash
}
\lstdefinestyle{python}{
    basicstyle=\ttfamily\footnotesize,
    backgroundcolor=\color{codebg},
    frame=single,
    rulecolor=\color{lightgray},
    breaklines=true,
    columns=fullflexible,
    keepspaces=true,
    showstringspaces=false,
    language=Python
}

\begin{document}

\begin{titlepage}
\thispagestyle{empty}
\vspace*{1.0cm}
\begin{center}
    {\Large\color{ineblue}\textbf{INSTITUTO NACIONAL ELECTORAL}}\\[0.8cm]
    {\Huge\color{ineblue}\textbf{Manual de usuario}}\\[0.25cm]
    {\LARGE Sistema de predicción y escenarios de riesgo de abstención}\\[1.0cm]
    \rule{0.85\textwidth}{0.8pt}\\[0.6cm]
    {\large Reto 2 — Conteos Censales de Participación Ciudadana}\\[0.25cm]
    {\large Versión del entregable: 2026}\\[1.2cm]
\end{center}

\begin{objetivo}
\textbf{Propósito del manual.} Este documento explica, paso a paso y sin asumir conocimientos de programación, cómo instalar, ejecutar, consultar y presentar el sistema. La persona lectora aprenderá a abrir el tablero, aplicar filtros territoriales, interpretar las gráficas, comparar escenarios y ubicar los archivos de resultados.
\end{objetivo}

\vfill
\begin{center}
\textbf{Equipo Loop-E}\\
\small Documento preparado para operación, demostración y reproducción del entregable.
\end{center}
\end{titlepage}

\tableofcontents
\newpage

\section{Resumen ejecutivo}

El sistema analiza los Conteos Censales de Participación Ciudadana (CCPC) del INE correspondientes al periodo 2009--2024. Su propósito es identificar territorios con mayor riesgo de abstención, estimar una proyección inercial hacia 2030 y comparar dicha proyección con un escenario en el que se aplican intervenciones de participación ciudadana.

El resultado principal para una persona usuaria no técnica es un tablero web autocontenido:

\begin{center}
\fbox{\parbox{0.9\textwidth}{
\centering
\ruta{outputs\textbackslash entregable\_escenarios\textbackslash mapa\_escenarios.html}
}}
\end{center}

El tablero permite consultar municipios y secciones, seleccionar una entidad federativa por nombre, filtrar municipios y secciones, observar gráficas dinámicas y comparar entidades o municipios entre sí.

\begin{resultado}
\textbf{Idea clave.} El tablero no pretende afirmar qué ocurrirá con certeza. Presenta escenarios analíticos para apoyar la priorización territorial y la planeación de pilotos institucionales.
\end{resultado}

\section{Qué puede hacer el sistema}

\subsection{Funciones principales}

\begin{itemize}
    \item Construir una base histórica de participación y abstención.
    \item Calcular un índice interpretable de riesgo territorial.
    \item Entrenar y validar un modelo de predicción temporal.
    \item Proyectar abstención a 2030 sin intervención.
    \item Simular una reducción hipotética de abstención con intervenciones.
    \item Generar gráficas nacionales, municipales y seccionales.
    \item Abrir un tablero interactivo sin instalar Power BI ni Tableau.
    \item Comparar entidades federativas, municipios y secciones.
    \item Exportar bases anonimizadas en formato CSV.
\end{itemize}

\subsection{Qué no hace el sistema}

\begin{itemize}
    \item No sustituye una decisión institucional.
    \item No identifica personas ni contiene datos personales individuales.
    \item No representa un efecto causal comprobado de las intervenciones.
    \item No es un mapa geográfico SHP, GeoJSON o GPKG; el entregable actual es un tablero territorial por códigos y nombres de consulta.
    \item No debe interpretarse una proyección como un resultado electoral observado.
\end{itemize}

\section{Conceptos explicados en lenguaje sencillo}

\subsection{Entidad federativa}

Es el estado o entidad de referencia, por ejemplo, Oaxaca, Chiapas o Nuevo León.

\subsection{Municipio}

Es una unidad territorial que agrupa secciones electorales. El tablero permite buscarlo por nombre.

\subsection{Sección electoral}

Es una unidad territorial más pequeña que el municipio. En el tablero se consulta mediante su número de sección.

\subsection{Abstención}

Es la proporción de personas de la lista nominal que no participó en una elección. En el tablero se muestra como porcentaje.

\subsection{Puntos porcentuales}

Una reducción de \textbf{3 puntos porcentuales} significa, por ejemplo, pasar de 50\% a 47\%. No significa reducir 3\% de 50\%.

\subsection{Riesgo territorial}

Es una clasificación analítica basada en tres señales:

\begin{enumerate}
    \item abstención reciente;
    \item tendencia histórica de participación;
    \item inestabilidad de la participación en los últimos ejercicios.
\end{enumerate}

Cada señal aporta un tercio al índice inicial. Las categorías son \alto{}, \medio{}, \bajo{} y \textit{sin histórico}.

\subsection{Escenario sin intervención}

Es una proyección inercial: supone que las tendencias históricas continúan sin aplicar una acción adicional.

\subsection{Escenario con intervención}

Es una simulación que aplica reducciones hipotéticas de abstención de acuerdo con el perfil de cada zona. Sus efectos deben validarse con pilotos.

\section{Requisitos para utilizar el programa}

\subsection{Equipo recomendado}

\begin{itemize}
    \item Windows 10/11, macOS o Linux.
    \item Al menos 8 GB de memoria RAM.
    \item 5 GB libres de almacenamiento.
    \item Navegador actualizado: Chrome, Edge, Firefox o Safari.
    \item Conexión a internet únicamente para descargar Python y dependencias; el tablero final funciona localmente.
\end{itemize}

\subsection{Programas necesarios}

\begin{enumerate}
    \item Python 3.10 o superior.
    \item Un editor de texto o Visual Studio Code.
    \item Un navegador web.
\end{enumerate}

\begin{importante}
\textbf{Para sólo consultar el tablero.} Si otra persona ya entregó la carpeta \ruta{outputs\textbackslash entregable\_escenarios}, no es necesario ejecutar el modelo completo. Basta con abrir \ruta{mapa\_escenarios.html}. La instalación de Python sólo es necesaria para reproducir o regenerar resultados.
\end{importante}

\section{Instalación para una persona sin experiencia técnica}

\subsection{Windows}

\begin{enumerate}
    \item Descargar Python desde \url{https://www.python.org/downloads/}.
    \item Ejecutar el instalador.
    \item Activar la casilla \textbf{Add Python to PATH}.
    \item Seleccionar \textbf{Install Now}.
    \item Abrir PowerShell desde el menú Inicio.
    \item Ubicar la carpeta del proyecto.
\end{enumerate}

Ejemplo de ubicación:

\begin{lstlisting}[style=terminal]
cd C:\Users\oscar\hackathon-ine
\end{lstlisting}

Verificar Python:

\begin{lstlisting}[style=terminal]
python --version
\end{lstlisting}

El resultado esperado es una versión como \comando{Python 3.10}, \comando{Python 3.11} o superior.

\subsection{macOS y Linux}

Abrir Terminal, ubicarse en la carpeta del repositorio y verificar:

\begin{lstlisting}[style=terminal]
cd /ruta/al/proyecto
python3 --version
\end{lstlisting}

En las instrucciones siguientes, sustituir \comando{python} por \comando{python3} si el sistema lo requiere.

\subsection{Instalar las dependencias}

Desde la raíz del proyecto:

\begin{lstlisting}[style=terminal]
python -m pip install -r requirements.txt
\end{lstlisting}

Las dependencias principales son:

\begin{itemize}
    \item \texttt{pandas}: lectura y transformación de datos.
    \item \texttt{numpy}: operaciones numéricas.
    \item \texttt{matplotlib}: generación de gráficas.
    \item \texttt{scikit-learn}: entrenamiento y validación del modelo.
\end{itemize}

\section{Estructura de carpetas}

\begin{longtable}{p{0.34\textwidth}p{0.58\textwidth}}
\toprule
\textbf{Carpeta} & \textbf{Uso}\\
\midrule
\ruta{data} & Datos originales, intermedios y procesados.\\
\ruta{src\textbackslash ingenieria\_datos} & Preparación y validación de la base analítica.\\
\ruta{src\textbackslash modelos} & Construcción de variables, riesgo, predicción y escenarios.\\
\ruta{src\textbackslash visualizaciones} & Gráficas y tablero interactivo.\\
\ruta{outputs\textbackslash analytical} & Series históricas agregadas.\\
\ruta{outputs\textbackslash features} & Variables de modelado.\\
\ruta{outputs\textbackslash risk} & Índices y niveles de riesgo.\\
\ruta{outputs\textbackslash escenarios} & Escenarios sin y con intervención.\\
\ruta{outputs\textbackslash visualizaciones} & Gráficas estáticas.\\
\ruta{outputs\textbackslash entregable\_escenarios} & Tablero web y CSV anonimizados.\\
\ruta{docs} & Informes y este manual.\\
\bottomrule
\end{longtable}

\section{Ejecución paso a paso del modelo}

\begin{importante}
\textbf{Orden obligatorio.} Los scripts deben ejecutarse en el orden indicado. Cada etapa utiliza archivos creados por la etapa anterior.
\end{importante}

\subsection{Paso 1: construir la base analítica}

\begin{lstlisting}[style=terminal]
python src\ingenieria_datos\03_construir_dataset_analitico.py
\end{lstlisting}

Esta etapa organiza la información histórica en dos niveles:

\begin{itemize}
    \item secciones electorales por año;
    \item municipios por año.
\end{itemize}

Archivos esperados:

\begin{lstlisting}[style=terminal]
outputs\analytical\ccpc_section_year.csv
outputs\analytical\ccpc_municipality_year.csv
\end{lstlisting}

\subsection{Paso 2: crear variables}

\begin{lstlisting}[style=terminal]
python src\modelos\01_creacion_variables_modelo.py
\end{lstlisting}

Se calculan variables como participación, abstención, cambios históricos, edad promedio, proporciones por grupo de edad y volatilidad.

\subsection{Paso 3: construir el índice de riesgo}

\begin{lstlisting}[style=terminal]
python src\modelos\02_construir_indice_riesgo.py
\end{lstlisting}

El resultado clasifica cada zona como \alto{}, \medio{}, \bajo{} o \textit{sin histórico}. Una zona sin histórico suficiente no debe interpretarse como una zona de bajo riesgo: significa que no existe evidencia temporal suficiente para calcular la tendencia.

\subsection{Paso 4: entrenar y validar}

\begin{lstlisting}[style=terminal]
python src\modelos\03_entrenar_y_validar_modelo.py
\end{lstlisting}

El modelo utiliza años anteriores a 2024 para entrenar y 2024 para validar. Los resultados se consultan en:

\begin{lstlisting}[style=terminal]
outputs\models\model_manifest.json
\end{lstlisting}

\subsection{Paso 5: integrar resultados}

\begin{lstlisting}[style=terminal]
python src\modelos\04_integrar_resultados_modelo.py
\end{lstlisting}

\textbf{Advertencia:} el archivo correcto contiene el nombre completo \texttt{04\_integrar\_resultados\_modelo.py}.

\subsection{Paso 6: generar rankings}

\begin{lstlisting}[style=terminal]
python src\modelos\05_generar_rankings_territoriales.py
\end{lstlisting}

\subsection{Paso 7: proyectar sin intervención}

\begin{lstlisting}[style=terminal]
python src\modelos\06_simular_escenario_sin_intervencion.py
\end{lstlisting}

Esta etapa produce la línea base de 2030:

\begin{itemize}
    \item \ruta{ccpc\_municipioEscenarioSinIntervencion.csv};
    \item \ruta{ccpc\_seccionEscenarioSinIntervencion.csv}.
\end{itemize}

\subsection{Paso 8: simular intervención}

\begin{lstlisting}[style=terminal]
python src\modelos\07_simular_escenario_intervencion.py
\end{lstlisting}

La simulación aplica tres propuestas:

\begin{description}
    \item[P1] Árbol de Influencers y Red de Embajadores Electorales.
    \item[P2] Campaña de Rescate Cívico con Monitoreo de Movilidad.
    \item[P3] Plataforma Cívica Digital Gamificada.
\end{description}

\section{Generación de gráficas}

\subsection{Panorama general}

\begin{lstlisting}[style=terminal]
python src\visualizaciones\01_generar_panorama_general.py
\end{lstlisting}

\subsection{Gráficas de intervención}

\begin{lstlisting}[style=terminal]
python src\visualizaciones\08_generar_visualizaciones_intervencion.py --solo-nacional
\end{lstlisting}

Las gráficas se guardan en:

\begin{lstlisting}[style=terminal]
outputs\visualizaciones\intervencion\municipios\
outputs\visualizaciones\intervencion\secciones\
\end{lstlisting}

\begin{longtable}{p{0.32\textwidth}p{0.60\textwidth}}
\toprule
\textbf{Archivo} & \textbf{Qué muestra}\\
\midrule
\texttt{01\_abstencionPorPerfil.png} & Abstención media sin y con intervención por perfil.\\
\texttt{02\_distribucionNiveles.png} & Cantidad de zonas por nivel de riesgo antes y después.\\
\texttt{03\_distribucionAbstencion.png} & Distribución de abstención proyectada.\\
\texttt{04\_baseVsIntervencion.png} & Comparación zona a zona; debajo de la diagonal hay mejora.\\
\texttt{05\_topZonasAntesDespues.png} & Diez zonas con mayor abstención proyectada.\\
\texttt{06\_reduccionPorEntidad.png} & Reducción media por entidad federativa.\\
\texttt{07\_sensibilidad.png} & Escenarios conservador, base y optimista.\\
\bottomrule
\end{longtable}

\section{Generar el tablero interactivo}

\subsection{Comando}

\begin{lstlisting}[style=terminal]
python src\visualizaciones\09_empaquetar_mapa_escenarios.py
\end{lstlisting}

\subsection{Archivos generados}

\begin{lstlisting}[style=terminal]
outputs\entregable_escenarios\mapa_escenarios.html
outputs\entregable_escenarios\municipios_escenarios_anonimizados.csv
outputs\entregable_escenarios\secciones_escenarios_anonimizados.csv
\end{lstlisting}

El HTML es autocontenido: las gráficas y los datos necesarios para la consulta quedan integrados en el archivo. No se necesita una base de datos, un backend ni conexión a internet.

\section{Cómo abrir el tablero}

\subsection{Método recomendado: servidor local}

Desde la raíz del proyecto ejecutar:

\begin{lstlisting}[style=terminal]
python -m http.server 5500
\end{lstlisting}

Abrir el navegador en:

\begin{lstlisting}[style=terminal]
http://127.0.0.1:5500/outputs/entregable_escenarios/mapa_escenarios.html
\end{lstlisting}

Para cerrar el servidor, regresar a la terminal y presionar \comando{Ctrl+C}.

\subsection{Visual Studio Code}

\begin{enumerate}
    \item Abrir la carpeta del proyecto.
    \item Abrir el archivo \ruta{outputs\textbackslash entregable\_escenarios\textbackslash mapa\_escenarios.html}.
    \item Hacer clic derecho.
    \item Seleccionar \textbf{Open with Live Server}.
    \item Si se abre la raíz del proyecto, agregar manualmente la ruta completa del archivo al navegador.
\end{enumerate}

\subsection{Abrir directamente}

También se puede hacer doble clic en el archivo HTML. Si el navegador bloquea alguna interacción, utilizar el servidor local de Python.

\section{Guía completa de uso del tablero}

\subsection{Pantalla inicial}

Al abrir el tablero se observan:

\begin{enumerate}
    \item encabezado y descripción del propósito;
    \item filtros territoriales;
    \item tarjetas de indicadores;
    \item gráficas comparativas;
    \item lectura ejecutiva;
    \item tabla de detalle.
\end{enumerate}

\subsection{Selector de alcance}

El selector \textbf{Alcance} tiene dos opciones:

\begin{itemize}
    \item \textbf{Municipios}: resume la información municipal.
    \item \textbf{Secciones}: permite llegar al nivel territorial más detallado.
\end{itemize}

\textbf{Procedimiento recomendado:}

\begin{enumerate}
    \item Elegir \textbf{Municipios} para una lectura general.
    \item Elegir una entidad federativa.
    \item Elegir un municipio.
    \item Revisar las tarjetas y gráficas.
    \item Cambiar a \textbf{Secciones} para investigar el interior del municipio.
\end{enumerate}

\subsection{Filtro por entidad federativa}

El selector muestra el nombre de la entidad. Al elegir una, se actualizan:

\begin{itemize}
    \item municipios disponibles;
    \item secciones disponibles;
    \item indicadores;
    \item gráficas;
    \item tabla de resultados.
\end{itemize}

\subsection{Filtro por municipio}

Después de elegir una entidad, el selector de municipio se llena con los municipios correspondientes. Al seleccionar un municipio, las métricas se recalculan sólo para esa unidad territorial.

\subsection{Filtro por sección}

En el alcance de secciones, el selector permite elegir la sección por número. Esta consulta sirve para observar el comportamiento de una zona específica.

\subsection{Filtro de riesgo}

El filtro de riesgo permite seleccionar:

\begin{itemize}
    \item \alto{};
    \item \medio{};
    \item \bajo{};
    \item \textit{sin histórico}.
\end{itemize}

\begin{importante}
\textbf{Sin histórico no significa bajo riesgo.} Significa que la zona no tiene suficientes observaciones para calcular una tendencia histórica confiable.
\end{importante}

\subsection{Restablecer filtros}

El botón \textbf{Restablecer filtros} devuelve el tablero a la vista inicial de municipios, todas las entidades, todos los municipios y todos los niveles de riesgo.

\section{Interpretación de las tarjetas}

\begin{longtable}{p{0.32\textwidth}p{0.58\textwidth}}
\toprule
\textbf{Tarjeta} & \textbf{Interpretación}\\
\midrule
Zonas con proyección & Número de zonas con información suficiente para proyectar.\\
Abstención sin intervención & Promedio del escenario inercial a 2030.\\
Abstención con intervención & Promedio después de aplicar los supuestos de intervención.\\
Reducción media & Diferencia promedio entre ambos escenarios, en puntos porcentuales.\\
Zonas en alto & Número de zonas clasificadas como alto antes y después de la simulación.\\
\bottomrule
\end{longtable}

\section{Interpretación de las gráficas}

\subsection{Abstención proyectada}

La gráfica compara dos barras:

\begin{itemize}
    \item rojo: escenario sin intervención;
    \item verde: escenario con intervención.
\end{itemize}

Una barra verde menor indica una reducción simulada de abstención. La reducción no debe interpretarse como una garantía.

\subsection{Nivel de riesgo}

La gráfica circular muestra cuántas zonas pertenecen a cada nivel de riesgo en el escenario base. El color rojo representa \alto{}, amarillo \medio{} y verde \bajo{}.

\subsection{Reducción por perfil}

Los perfiles son:

\begin{description}
    \item[A, base crítica] zonas con abstención base elevada o riesgo alto.
    \item[B, deterioro] zonas con tendencia descendente de participación.
    \item[C, inestable] zonas con señal de inestabilidad cíclica.
    \item[Sin perfil] zonas sin una señal específica de intervención.
\end{description}

\subsection{Distribución de abstención}

Esta vista agrupa las zonas en rangos de abstención. Permite observar si, después de la intervención simulada, se desplaza la distribución hacia valores menores.

\subsection{Comparación por entidad}

La lista ordena entidades por reducción media de abstención. También muestra el cambio en el número de zonas clasificadas como \alto{}.

\subsection{Prioridad territorial}

El ranking muestra las diez zonas con mayor abstención proyectada sin intervención. Debe utilizarse como punto de partida para revisar contexto territorial, no como una orden automática de intervención.

\section{Comparador territorial}

\subsection{Objetivo}

El comparador permite responder preguntas como:

\begin{itemize}
    \item ¿Qué entidad tiene mayor abstención proyectada?
    \item ¿Cuál entidad presenta mayor reducción simulada?
    \item ¿Cómo se compara un municipio con otro?
    \item ¿Qué municipio mantiene más zonas en riesgo alto?
\end{itemize}

\subsection{Procedimiento}

\begin{enumerate}
    \item Seleccionar una entidad en \textbf{Entidad A}.
    \item Seleccionar otra entidad en \textbf{Entidad B}.
    \item Opcionalmente seleccionar un municipio en \textbf{Municipio A}.
    \item Opcionalmente seleccionar otro municipio en \textbf{Municipio B}.
    \item Observar las barras rojas y verdes.
\end{enumerate}

La barra roja corresponde al escenario sin intervención y la verde al escenario con intervención. El número de zonas indica cuántas observaciones forman el promedio.

\begin{resultado}
\textbf{Buena práctica.} Comparar territorios con el mismo alcance. No comparar directamente un municipio con una sección sin aclarar que representan unidades territoriales de distinta escala.
\end{resultado}

\section{Tabla de detalle}

La tabla contiene:

\begin{itemize}
    \item identificador anonimizado de zona;
    \item código de entidad;
    \item código de municipio;
    \item abstención sin intervención;
    \item abstención con intervención;
    \item reducción en puntos porcentuales;
    \item nivel de riesgo base;
    \item nivel de riesgo después de la intervención;
    \item propuesta aplicada.
\end{itemize}

La tabla presenta hasta 500 registros visibles por filtro. Para análisis masivo, utilizar los CSV anonimizados.

\section{Bases anonimizadas}

\subsection{Municipios}

\ruta{outputs\textbackslash entregable\_escenarios\textbackslash municipios\_escenarios\_anonimizados.csv}

\subsection{Secciones}

\ruta{outputs\textbackslash entregable\_escenarios\textbackslash secciones\_escenarios\_anonimizados.csv}

\subsection{Campos principales}

\begin{longtable}{p{0.39\textwidth}p{0.50\textwidth}}
\toprule
\textbf{Campo} & \textbf{Descripción}\\
\midrule
\texttt{zona\_id} & Identificador hash anonimizado.\\
\texttt{nivel\_territorial} & Municipio o sección.\\
\texttt{entidad\_codigo} & Código de entidad.\\
\texttt{municipio\_codigo} & Código de municipio.\\
\texttt{seccion\_codigo} & Código de sección, cuando corresponde.\\
\texttt{abstencion\_sin\_intervencion\_pct} & Abstención proyectada sin intervención.\\
\texttt{abstencion\_con\_intervencion\_pct} & Abstención proyectada con intervención.\\
\texttt{reduccion\_abstencion\_pp} & Reducción simulada en puntos porcentuales.\\
\texttt{nivel\_riesgo\_base} & Nivel de riesgo antes de intervención.\\
\texttt{nivel\_riesgo\_intervencion} & Nivel resultante bajo el escenario simulado.\\
\texttt{perfil\_intervencion} & Perfil A, B, C o sin perfil.\\
\texttt{propuestas\_aplicadas} & Propuesta o combinación aplicada.\\
\bottomrule
\end{longtable}

\section{Flujo completo reproducible}

El siguiente bloque reproduce todo el pipeline en Windows PowerShell:

\begin{lstlisting}[style=terminal]
cd C:\Users\oscar\hackathon-ine
python -m pip install -r requirements.txt
python src\ingenieria_datos\03_construir_dataset_analitico.py
python src\modelos\01_creacion_variables_modelo.py
python src\modelos\02_construir_indice_riesgo.py
python src\modelos\03_entrenar_y_validar_modelo.py
python src\modelos\04_integrar_resultados_modelo.py
python src\modelos\05_generar_rankings_territoriales.py
python src\modelos\06_simular_escenario_sin_intervencion.py
python src\modelos\07_simular_escenario_intervencion.py
python src\visualizaciones\01_generar_panorama_general.py
python src\visualizaciones\08_generar_visualizaciones_intervencion.py --solo-nacional
python src\visualizaciones\09_empaquetar_mapa_escenarios.py
python -m http.server 5500
\end{lstlisting}

Después abrir:

\begin{lstlisting}[style=terminal]
http://127.0.0.1:5500/outputs/entregable_escenarios/mapa_escenarios.html
\end{lstlisting}

\section{Solución de problemas}

\subsection{El comando \texttt{python} no existe}

Instalar Python y activar \textbf{Add Python to PATH}. Cerrar y volver a abrir la terminal.

\subsection{No se encuentra un archivo CSV}

Verificar que los scripts se ejecutaron desde la raíz del repositorio y en orden. No saltar los pasos de construcción de la base, integración y escenarios.

\subsection{La página muestra la carpeta del proyecto}

Se abrió el servidor en la raíz, pero no la página. Agregar la ruta:

\begin{lstlisting}[style=terminal]
http://127.0.0.1:5500/outputs/entregable_escenarios/mapa_escenarios.html
\end{lstlisting}

\subsection{La página aparece en blanco}

\begin{enumerate}
    \item Regenerar el tablero:
\begin{lstlisting}[style=terminal]
python src\visualizaciones\09_empaquetar_mapa_escenarios.py
\end{lstlisting}
    \item Actualizar el navegador con \comando{Ctrl+F5}.
    \item Utilizar \comando{python -m http.server 5500}.
    \item Revisar que exista \ruta{mapa\_escenarios.html}.
\end{enumerate}

\subsection{El selector de municipios está vacío}

Primero seleccionar una entidad. Si sigue vacío, cambiar a \textbf{Municipios}, restablecer filtros y regenerar el tablero.

\subsection{El modelo tarda}

El procesamiento puede tardar debido al volumen histórico. No cerrar la terminal mientras el script esté trabajando. Esperar a que aparezca el mensaje final del proceso.

\subsection{Aparece una advertencia de \textit{sin histórico}}

No es un error. Indica que no hay suficientes años observados para construir una tendencia confiable.

\section{Validación antes de entregar}

\begin{enumerate}
    \item Confirmar que existe \ruta{mapa\_escenarios.html}.
    \item Abrir el tablero en el navegador.
    \item Probar el filtro de entidad.
    \item Probar el filtro de municipio.
    \item Probar una sección.
    \item Probar el botón de restablecer filtros.
    \item Ejecutar una comparación entidad contra entidad.
    \item Ejecutar una comparación municipio contra municipio.
    \item Confirmar que las gráficas cambian al aplicar filtros.
    \item Confirmar que los dos CSV anonimizados existen.
\end{enumerate}

\section{Resultados de referencia}

\begin{longtable}{p{0.46\textwidth}r}
\toprule
\textbf{Indicador} & \textbf{Valor}\\
\midrule
Secciones clasificadas & 70,751\\
Municipios clasificados & 2,477\\
MAE Random Forest, secciones & 0.0775\\
MAE baseline, secciones & 0.0849\\
MAE Random Forest, municipios & 0.0700\\
MAE baseline, municipios & 0.0660\\
\bottomrule
\end{longtable}

\section{Limitaciones y uso responsable}

\begin{itemize}
    \item Los supuestos de P1, P2 y P3 son hipótesis de diseño.
    \item La simulación no estima causalidad.
    \item Las zonas sin histórico requieren una estrategia separada de fortalecimiento de datos.
    \item Los resultados deben revisarse con especialistas territoriales.
    \item La asignación de recursos debe realizarse después de validar los efectos mediante pilotos.
\end{itemize}

\section{Glosario rápido}

\begin{description}
    \item[Abstención] Porcentaje estimado de no participación.
    \item[CCPC] Conteos Censales de Participación Ciudadana.
    \item[CSV] Archivo tabular separado por comas.
    \item[Escenario base] Proyección sin intervención.
    \item[Escenario con intervención] Proyección después de aplicar supuestos de reducción.
    \item[Índice de riesgo] Puntaje compuesto de abstención, tendencia e inestabilidad.
    \item[MAE] Error absoluto medio de una predicción.
    \item[PP] Puntos porcentuales.
    \item[Proyección] Estimación futura basada en información histórica.
\end{description}

\section{Checklist para una demostración}

\begin{enumerate}
    \item Abrir la URL local del tablero.
    \item Explicar que el rojo es sin intervención y el verde con intervención.
    \item Mostrar el total de zonas con proyección.
    \item Seleccionar una entidad federativa.
    \item Seleccionar un municipio.
    \item Cambiar al alcance de secciones.
    \item Mostrar una sección específica.
    \item Comparar dos entidades.
    \item Comparar dos municipios.
    \item Explicar que los resultados son escenarios y no garantías.
\end{enumerate}

\begin{resultado}
\textbf{Mensaje final para la presentación.} El sistema transforma datos históricos en una herramienta de priorización territorial: muestra dónde puede concentrarse el riesgo, cómo cambiaría bajo distintos supuestos y qué información debe validarse antes de intervenir.
\end{resultado}

\end{document}
